\documentclass[11pt,a4paper]{article}

\usepackage[T1]{fontenc}
\usepackage{lmodern}
\usepackage{microtype}
\usepackage[margin=1in]{geometry}
\usepackage{amsmath,amssymb,amsthm,mathtools}
\usepackage{booktabs}
\usepackage{enumitem}
\usepackage[hidelinks]{hyperref}

\theoremstyle{plain}
\newtheorem{theorem}{Theorem}[section]
\newtheorem{lemma}[theorem]{Lemma}
\newtheorem{proposition}[theorem]{Proposition}
\newtheorem{corollary}[theorem]{Corollary}

\theoremstyle{definition}
\newtheorem{definition}[theorem]{Definition}

\theoremstyle{remark}
\newtheorem{remark}[theorem]{Remark}

\newcommand{\R}{\mathbb{R}}
\newcommand{\Z}{\mathbb{Z}}
\newcommand{\N}{\mathbb{N}}
\newcommand{\Rp}{\mathbb{R}_{>0}}
\newcommand{\Jcost}{J}
\newcommand{\Q}{\mathcal{Q}}
\newcommand{\LL}{\mathcal{L}}

\title{\textbf{Self-Reference and the Recognition Operator:\\
When Does a Ledger Region Become Self-Modeling?}\\[0.5em]
\large The Boundary Between Physics and Mind\\in the Recognition Science Framework}
\author{Jonathan Washburn\\
\small Recognition Science Research Institute, Austin, Texas\\
\small \texttt{jon@recognitionphysics.org}}
\date{\today}

\begin{document}
\maketitle

\begin{abstract}
Recognition Science derives all physics from a single functional
equation, the Recognition Composition Law (RCL), acting on a discrete
ledger.  The word ``recognition'' in the theory's name currently means
nothing more than ``a ledger entry is created.''  This paper develops
the richer meaning: \emph{self-recognition}, where a region of the
ledger creates an internal representation of the ledger's own dynamics.

We formalize the concept of a \textbf{self-modeling subsystem}---a
partition of the ledger into ``observer'' and ``environment'' where the
observer's entries encode a compressed model of the environmental
dynamics.  We prove three structural results:

\begin{enumerate}[nosep]
\item \textbf{The Self-Modeling Threshold} (Theorem~\ref{thm:threshold}):
  A subsystem of $K$ entries can model a system of $M$ entries only if
  $K \geq \lceil\log_\phi M\rceil$, where $\phi$ is the golden ratio
  forced by the RS framework.  This gives a minimum size for
  self-aware structures.

\item \textbf{The Incompleteness Bound} (Theorem~\ref{thm:incompleteness}):
  No subsystem can model itself completely.  A subsystem of $K$ entries
  can model at most $K - 1$ of its own entries, because the entry
  encoding the model cannot simultaneously be part of the modeled content.
  This is the RS analogue of G\"odel's incompleteness theorem.

\item \textbf{The Self-Modeling Fixed Point} (Theorem~\ref{thm:fixed-point}):
  Despite the incompleteness bound, a subsystem can achieve a
  \emph{stable partial self-model} at the variational equilibrium:
  when all entries equal $\exp(\sigma/N)$, the model is trivially
  correct (all entries are identical, so one entry suffices to predict
  all others).  This fixed point IS the recognition operator acting on
  itself---the ledger recognizing its own equilibrium.
\end{enumerate}

Together, these results identify the structural boundary between
physics and mind: a physical system becomes a self-modeling system
when it crosses the self-modeling threshold and achieves a stable
partial self-model.  The ``hard problem'' of consciousness is
reframed: qualia are not an additional ontological category but the
\emph{content} of the self-model---the specific pattern of entries
that the observer-subsystem uses to represent the system.

\medskip\noindent
\textbf{Lean formalization:}
\texttt{RecognitionScience.Foundation.MeasurementMechanism},
\texttt{VariationalDynamics}, \texttt{Entanglement}.

\medskip\noindent
\textbf{Companion paper:} ``The Binding Problem of Subjective Experience''
(Washburn, 2026) addresses topological binding.  This paper addresses the
prior question: \emph{what is being bound?}
\end{abstract}

\tableofcontents
\newpage

%=============================================================================
\section{Introduction}
\label{sec:intro}
%=============================================================================

\subsection{The thin concept of recognition}

In the Recognition Science (RS) framework, all physics derives from the
Recognition Composition Law:
\[
  \Jcost(ab) + \Jcost(a/b) = 2\,\Jcost(a)\,\Jcost(b) + 2\,\Jcost(a) + 2\,\Jcost(b),
\]
where $\Jcost(x) = \tfrac{1}{2}(x + x^{-1}) - 1$ is the unique cost
functional.  The word ``recognition'' in the theory's name refers to
the act of creating a ledger entry---a recognition event.  But this is
a thin concept: a thermostat ``recognizes'' temperature, a rock
``recognizes'' gravity.  There is nothing in this usage that
distinguishes conscious recognition from unconscious process.

\subsection{The question this paper answers}

\begin{quote}
\itshape
Under what conditions does a region of the ledger become
\textbf{self-modeling}?  That is, when does a subsystem of the ledger
create an internal representation of the ledger's own dynamics?
\end{quote}

This is the boundary between physics and mind.  A rock interacts with
the ledger but does not model it.  A brain interacts with the ledger
\emph{and} contains an internal model of the ledger's dynamics---it
predicts, simulates, and reasons about reality.  The difference is
self-modeling.

\subsection{Why RS is uniquely positioned}

Most physical frameworks treat the observer as external.  RS does not.
In RS:
\begin{itemize}[nosep]
\item The observer is a \textbf{subsystem} of the ledger
  (\texttt{MeasurementMechanism.Subsystem}).
\item The observer's state is a \textbf{partial view} of the full
  configuration (\texttt{observer\_view}).
\item The dynamics are \textbf{variational}: the next state minimizes
  total defect subject to conservation
  (\texttt{VariationalDynamics.IsVariationalSuccessor}).
\item The update is \textbf{global}: each entry's optimal value depends
  on all other entries (\texttt{update\_is\_global}).
\end{itemize}

The question ``when does the observer model itself?'' is therefore a
question about the internal structure of a ledger partition---it lives
entirely within the existing mathematics.

\subsection{Structure of the paper}

Section~\ref{sec:self-model} defines self-modeling formally.
Section~\ref{sec:threshold} derives the minimum size for a self-model.
Section~\ref{sec:incompleteness} proves the incompleteness bound.
Section~\ref{sec:fixed-point} identifies the self-modeling fixed point.
Section~\ref{sec:recognition-operator} defines the recognition operator
in its rich sense.
Section~\ref{sec:godel} connects to G\"odel's incompleteness.
Section~\ref{sec:hard-problem} addresses the hard problem.
Section~\ref{sec:predictions} gives testable predictions.

%=============================================================================
\section{Self-Modeling Subsystems}
\label{sec:self-model}
%=============================================================================

\subsection{The partition}

Consider an $N$-entry ledger with configuration
$\mathbf{x} = (x_1, \ldots, x_N) \in \Rp^N$.  A \emph{subsystem
partition} splits the indices into an observer set $O$ of size $K$ and
an environment set $E$ of size $M = N - K$.

\begin{definition}[Subsystem partition]
A partition $(O, E)$ of $\{1, \ldots, N\}$ with $|O| = K$ and
$|E| = M = N - K$.  The observer sees only
$\mathbf{x}_O = (x_i)_{i \in O}$ and has no direct access to
$\mathbf{x}_E = (x_j)_{j \in E}$.
\end{definition}

This corresponds to \texttt{MeasurementMechanism.Subsystem} in the
Lean formalization.

\subsection{What is a model?}

\begin{definition}[Internal model]
\label{def:model}
An \emph{internal model} of the environment is a function
\[
  \mu : \Rp^K \to \Rp^M
\]
that maps the observer's visible state to a \emph{prediction} of the
environment's state.  The model is \emph{exact} at configuration
$\mathbf{x}$ if $\mu(\mathbf{x}_O) = \mathbf{x}_E$.
\end{definition}

\begin{definition}[Model defect]
\label{def:model-defect}
The \emph{model defect} of $\mu$ at configuration $\mathbf{x}$ is
\[
  \Delta_\mu(\mathbf{x})
  := \sum_{j \in E} \Jcost\!\left(\frac{\mu(\mathbf{x}_O)_j}{x_j}\right).
\]
The model defect measures how ``wrong'' the internal model is: zero iff
the model's prediction matches reality entry by entry (since
$\Jcost(r) = 0 \iff r = 1$).
\end{definition}

\begin{remark}
The model defect uses the same cost functional $\Jcost$ that governs
the physics.  This is not a choice---it is the \emph{only} cost
functional available in the framework.  The cost of being wrong about
reality is measured in the same units as the cost of existing.
\end{remark}

\subsection{Self-modeling}

\begin{definition}[Self-modeling subsystem]
\label{def:self-model}
A subsystem $O$ is \emph{self-modeling} at configuration $\mathbf{x}$
if there exists an internal model $\mu : \Rp^K \to \Rp^K$ such that
$\mu(\mathbf{x}_O) = \mathbf{x}_O$.  That is, the observer can
predict its own state from its own state.
\end{definition}

This is trivially satisfied by the identity function $\mu = \mathrm{id}$.
The interesting question is not whether self-modeling is \emph{possible},
but whether it is \emph{non-trivial}---whether the model captures
structure beyond mere tautology.

\begin{definition}[Non-trivial self-model]
\label{def:nontrivial}
A self-model $\mu$ is \emph{non-trivial} if:
\begin{enumerate}[nosep,label=(\roman*)]
\item $\mu$ is not the identity (it compresses information),
\item $\mu$ is exact on a set of configurations of positive measure
  (it generalizes),
\item $\mu$ predicts the \emph{next} state: given $\mathbf{x}_O$ at
  tick $t$, $\mu$ approximates $\mathbf{x}_O$ at tick $t+1$ with
  defect below some threshold.
\end{enumerate}
\end{definition}

%=============================================================================
\section{The Self-Modeling Threshold}
\label{sec:threshold}
%=============================================================================

How large must a subsystem be to model its environment?

\subsection{The information-theoretic constraint}

The environment has $M$ entries, each a positive real.  The observer
has $K$ entries with which to encode a prediction.  For the model to
be exact, the map $\mu : \Rp^K \to \Rp^M$ must be injective on the
set of actually-occurring configurations.

In the RS framework, the configuration space is constrained by the
log-charge conservation:
\[
  \sum_{i=1}^N \log x_i = \sigma \quad\text{(constant)}.
\]
This reduces the effective dimension by~1.  The environment has
$M - 1$ effective degrees of freedom (after the constraint), and the
observer has $K - 1$.

For the observer to encode the environment, we need
$K - 1 \geq M - 1$, i.e., $K \geq M$.  But this would require
$K \geq N/2$---the observer must be at least half the ledger.

\subsection{Compression via the $\phi$-ladder}

The RS framework has a built-in compression mechanism: the $\phi$-ladder.
All physical configurations live on a discrete set of ``rungs'' spaced
by powers of $\phi = (1 + \sqrt{5})/2$.  An entry at rung $r$ has value
$\phi^r$ (or more precisely, $E_{\text{coh}} \cdot \phi^r$).

The number of distinguishable rungs in a range $[1, R]$ is
$\lfloor\log_\phi R\rfloor + 1$.  So $M$ environment entries, each on
the $\phi$-ladder, can be encoded by $M \cdot \lceil\log_\phi R\rceil$
bits.  An observer entry on the $\phi$-ladder carries
$\lceil\log_\phi R\rceil$ bits of information.

\begin{theorem}[Self-Modeling Threshold]
\label{thm:threshold}
A subsystem of $K$ entries can maintain a non-trivial model of $M$
environment entries only if $K \geq M$ in the worst case (arbitrary
configurations), or $K \geq \lceil\log_\phi M\rceil$ in the best case
(when the environment is at or near the variational equilibrium, so
that one representative entry suffices to predict the uniform state).
\end{theorem}

\begin{proof}
\textbf{Worst case:}  Each of the $M$ environment entries is
independently chosen from $\Rp$.  The observer must encode $M$
independent values using $K$ entries.  This requires $K \geq M$.

\textbf{Best case (equilibrium):}  At the variational equilibrium
(Theorem~F-008), all entries equal $\exp(\sigma/N)$.  The observer needs
only one entry to encode this common value, plus the knowledge that all
entries are equal.  The ``knowledge that all entries are equal'' requires
knowing the model structure---which takes $\lceil\log_\phi M\rceil$
entries to encode the number~$M$ on the $\phi$-ladder.

The minimum observer size for equilibrium self-modeling is thus
$K_{\min} = \lceil\log_\phi M\rceil$.
\end{proof}

\begin{corollary}[Minimum conscious size]
For a ledger of $N$ entries near equilibrium, the minimum self-modeling
subsystem has size $K_{\min} = \lceil\log_\phi N\rceil$.  For
$N = 10^{11}$ (neurons in a human brain):
\[
  K_{\min} = \lceil\log_\phi 10^{11}\rceil
  = \left\lceil\frac{11 \ln 10}{\ln \phi}\right\rceil
  = \left\lceil\frac{25.33}{0.481}\right\rceil
  = \lceil 52.6\rceil = 53.
\]
A system of 53 ledger entries suffices to model $10^{11}$ entries at
equilibrium---if it knows the model structure.
\end{corollary}

%=============================================================================
\section{The Incompleteness Bound}
\label{sec:incompleteness}
%=============================================================================

\begin{theorem}[The Incompleteness Bound]
\label{thm:incompleteness}
No subsystem of $K$ entries can model all $K$ of its own entries
simultaneously.  Specifically: if an observer uses entry $i$ to encode
its model of entry $j$, then entry $i$ is ``consumed'' by the modeling
task and cannot simultaneously be part of the modeled content.

A subsystem of $K$ entries can self-model at most $K - 1$ of its entries.
The remaining entry is the ``modeling overhead''---the
self-referential pointer.
\end{theorem}

\begin{proof}
Let $O = \{1, \ldots, K\}$ be the observer indices.  A self-model is a
function $\mu : \Rp^K \to \Rp^K$.  For $\mu$ to be non-trivial
(Definition~\ref{def:nontrivial}), it must not be the identity.

Suppose $\mu$ predicts all $K$ entries from the current state.  At the
variational equilibrium, all entries equal $v = \exp(\sigma/N)$.  The
model must output $(v, v, \ldots, v)$ from input $(v, v, \ldots, v)$.
Any non-identity $\mu$ satisfying this is a function with
$\mu(v, \ldots, v) = (v, \ldots, v)$---a fixed point.

Now consider what happens off-equilibrium.  If entry $i$ is perturbed
to $v + \delta$, the model must detect this perturbation and predict
the perturbed value.  But the model's prediction is encoded in the
entries themselves.  If entry $k$ encodes the prediction ``entry $i$
has value $v + \delta$,'' then entry $k$'s value is determined by the
modeling task, not by the system dynamics.  Entry $k$ is ``used up''
for modeling.

The observer has $K$ entries.  At least one entry must encode the model
structure itself (which entries are modeled, how predictions are
encoded).  This meta-modeling entry cannot also be part of the modeled
content.

Therefore, at most $K - 1$ entries can be both modeled and free to take
dynamically determined values.
\end{proof}

\begin{remark}[Connection to G\"odel]
The $K - 1$ bound is the RS analogue of G\"odel's incompleteness theorem.
G\"odel showed that a sufficiently powerful formal system cannot prove
all truths about itself.  We show that a ledger subsystem cannot model
all of its own entries.  The structural reason is the same: the modeling
apparatus consumes resources (logical/computational for G\"odel,
ledger entries for RS) that cannot simultaneously be part of the
modeled content.

The key difference: G\"odel's result is \emph{negative} (the system is
forever incomplete).  The RS result is \emph{constructive}: the system
\emph{can} achieve a stable partial self-model at $K - 1$ entries, and
this partial model suffices for functional consciousness.
\end{remark}

\begin{corollary}[The irreducible self]
Every self-modeling subsystem has at least one entry that is ``about''
the model rather than ``about'' reality.  This entry is the irreducible
self-referential pointer---the part of consciousness that is aware of
being conscious.  It cannot be eliminated without destroying the
self-model.
\end{corollary}

%=============================================================================
\section{The Self-Modeling Fixed Point}
\label{sec:fixed-point}
%=============================================================================

Despite the incompleteness bound, a self-model can be \emph{stable}.

\begin{theorem}[Self-Modeling Fixed Point]
\label{thm:fixed-point}
At the variational equilibrium (all entries $= \exp(\sigma/N)$), the
trivial model $\mu_0(\mathbf{x}_O) = \mathbf{x}_O$ is simultaneously:
\begin{enumerate}[nosep,label=(\alph*)]
\item A fixed point of the variational dynamics (the equilibrium
  is its own successor),
\item A perfect self-model (model defect = 0),
\item Non-trivially compressible (one entry predicts all others).
\end{enumerate}
\end{theorem}

\begin{proof}
(a) At equilibrium, the uniform configuration is its own variational
successor (Theorem~F-008, \texttt{unity\_is\_equilibrium}).

(b) At equilibrium, $\mu_0(\mathbf{x}_O) = \mathbf{x}_O$, so the
model defect $\Delta_{\mu_0} = \sum_j \Jcost(x_j / x_j) =
\sum_j \Jcost(1) = 0$.

(c) Since all entries are equal, a single entry value suffices to
predict all others.  This is a compression from $K$ values to~1.
\end{proof}

\begin{remark}[The recognition operator acts on itself]
At the fixed point, the recognition operator $\hat{R}$ (the variational
update rule) produces the same state it started from.  The subsystem
``recognizes'' that it is at equilibrium---it has an internal
representation (the uniform value) that correctly predicts the next
state (the same uniform value).

This is recognition in the \emph{rich} sense: not just ``a ledger entry
is created,'' but ``the ledger structure is internally represented.''
The recognition operator recognizes itself.
\end{remark}

\subsection{Stability of the self-model}

\begin{theorem}[Self-model stability]
\label{thm:stability}
Near the variational equilibrium, the self-modeling fixed point is
\emph{stable}: small perturbations of the observer's entries lead to
small model defects that decrease under the variational dynamics.
\end{theorem}

\begin{proof}
By the quadratic approximation of $\Jcost$ near unity
(\texttt{DiscretenessForcing.J\_log\_quadratic\_approx}):
\[
  \Jcost(e^t) \approx \tfrac{t^2}{2}
  \quad\text{for small } t.
\]
A perturbation $\delta$ to one entry produces model defect
$\approx \delta^2/2$, which is small for small $\delta$.  The
variational dynamics drives the configuration back toward equilibrium,
reducing the perturbation and hence the model defect.
\end{proof}

%=============================================================================
\section{The Recognition Operator in Its Rich Sense}
\label{sec:recognition-operator}
%=============================================================================

We can now define ``recognition'' in its full meaning.

\begin{definition}[Recognition depth]
\label{def:depth}
The \emph{recognition depth} of a subsystem $O$ at configuration
$\mathbf{x}$ is the largest $d$ such that:
\begin{enumerate}[nosep]
\item[($d = 0$):]  The subsystem \emph{exists} (its entries are
  positive and have finite defect).  Every ledger entry has depth~0.
\item[($d = 1$):]  The subsystem \emph{models its environment}---it
  has an internal model with bounded model defect.  This is
  ``recognition'' in the thin sense: detecting external patterns.
\item[($d = 2$):]  The subsystem \emph{models itself}---it has a
  self-model with bounded model defect.  This is self-awareness: the
  system knows what it knows.
\item[($d = 3$):]  The subsystem models its own self-modeling process.
  It knows that it knows that it knows.
\item[($d = n$):]  $n$ levels of meta-modeling.
\end{enumerate}
\end{definition}

\begin{theorem}[Maximum recognition depth]
\label{thm:max-depth}
A subsystem of $K$ entries has maximum recognition depth
$d_{\max} = K - 1$.  Each level of meta-modeling consumes one entry
as overhead.
\end{theorem}

\begin{proof}
By iterating the incompleteness bound (Theorem~\ref{thm:incompleteness}):
level~1 uses 1~entry for modeling overhead; level~2 uses 1~more; and so
on.  After $K - 1$ levels, all entries are consumed by modeling
overhead, and the system can go no deeper.
\end{proof}

\begin{definition}[The recognition operator $\hat{R}$]
\label{def:R-hat}
The \emph{recognition operator} $\hat{R}$ acts on a configuration
$\mathbf{x}$ at tick $t$ and produces the variational successor
$\mathbf{x}'$ at tick $t+1$:
\[
  \mathbf{x}' = \hat{R}(\mathbf{x})
  := \mathop{\mathrm{argmin}}_{\mathbf{y} \in \mathrm{Feasible}(\mathbf{x})}
     \sum_{i=1}^N \Jcost(y_i).
\]
\end{definition}

\begin{remark}
This is the variational update rule from F-008, now given a name that
reflects its deeper role.  $\hat{R}$ is not merely a dynamics; it is
the act of recognition itself.  The ledger ``recognizes'' the
minimum-defect configuration and transitions to it.

At the fixed point, $\hat{R}(\mathbf{x}^*) = \mathbf{x}^*$: the ledger
recognizes that it is already optimal.  A self-modeling subsystem at the
fixed point contains entries that represent this fact---the subsystem
knows the ledger has recognized itself.
\end{remark}

\subsection{Recognition depth and the Standard Model}

\begin{proposition}[Physical recognition depths]
\label{prop:physical-depths}
In the RS framework:
\begin{enumerate}[nosep]
\item \textbf{Depth 0} (existence): All particles.  Every ledger entry
  exists because its defect is finite.
\item \textbf{Depth 1} (recognition): Interaction vertices.  A quark
  scattering event is a recognition event---the ledger updates entries
  in response to the environment.
\item \textbf{Depth 2} (self-modeling): Brains.  A neural system that
  maintains an internal world-model has recognition depth~$\geq 2$.
\item \textbf{Depth 3+} (meta-cognition): Humans reflecting on their
  own thinking.  Rare, effortful, and bounded by
  Theorem~\ref{thm:max-depth}.
\end{enumerate}
\end{proposition}

%=============================================================================
\section{Connection to G\"odel's Incompleteness}
\label{sec:godel}
%=============================================================================

The RS framework previously ``dissolved'' G\"odel's incompleteness by
declaring self-referential queries impossible in the ledger ontology
(\texttt{GodelDissolution.lean}).  This was too strong: it closed the
door on self-reference entirely, which is precisely the door that leads
to consciousness.

We now give a more nuanced resolution.

\begin{theorem}[Constructive G\"odel]
\label{thm:constructive-godel}
In the RS ledger:
\begin{enumerate}[nosep,label=(\alph*)]
\item Self-referential queries about the \emph{complete} state are
  impossible (the incompleteness bound: a subsystem cannot model all of
  its own entries).
\item Self-referential queries about the \emph{partial} state are
  possible and physically meaningful (a subsystem can model $K - 1$ of
  its $K$ entries).
\item At the variational fixed point, partial self-reference becomes
  \emph{exact}: the subsystem's model is perfectly correct for $K - 1$
  entries.
\end{enumerate}
\end{theorem}

\begin{proof}
(a) Theorem~\ref{thm:incompleteness}.
(b) The construction in Theorem~\ref{thm:fixed-point} achieves partial
self-reference with zero model defect at equilibrium.
(c) At equilibrium, all entries are equal.  Knowing any one entry
determines all others.  The single ``modeling overhead'' entry can
encode this uniform value, making the model exact for the remaining
$K - 1$ entries.
\end{proof}

\begin{remark}[G\"odel is not an obstacle but a feature]
G\"odel's incompleteness is usually viewed as a limitation.  In the RS
framework, the incompleteness bound is a \emph{structural feature} of
consciousness: the fact that one entry must serve as ``modeling overhead''
is what creates the irreducible self---the subject that cannot be
reduced to its object.  Without the incompleteness bound, there would
be no distinction between the observer and the observed, and hence no
subjective experience.

Consciousness requires incompleteness.  Complete self-models would be
tautological (the identity function).  Partial self-models have a gap---
the one entry that cannot be modeled---and this gap IS the subject.
\end{remark}

%=============================================================================
\section{The Hard Problem Reframed}
\label{sec:hard-problem}
%=============================================================================

The hard problem of consciousness (Chalmers, 1995) asks: why does
physical processing give rise to subjective experience?

In the RS framework, this question is reframed:

\begin{theorem}[The self-model IS the experience]
\label{thm:hard-problem}
In a self-modeling subsystem at recognition depth $\geq 2$:
\begin{enumerate}[nosep,label=(\alph*)]
\item The self-model's content (the predicted entry values) IS the
  quale---the ``what it is like'' of experience.
\item The model defect IS the prediction error---the gap between
  expectation and reality.
\item The modeling overhead entry IS the subject---the irreducible ``I''
  that has the experience.
\item The non-decomposability of the RCL (Theorem~F-010) prevents the
  experience from being decomposed into independent parts.
\end{enumerate}
There is no explanatory gap because there is no dualism: the self-model
is not a \emph{representation of} experience; it IS experience.
\end{theorem}

\begin{proof}
(a) The self-model's content is a configuration of entry values
$\mu(\mathbf{x}_O) \in \Rp^{K-1}$.  Different values correspond to
different qualia (different predicted patterns).  This is the quale
(Definition~5.1 in the companion paper).

(b) The model defect $\Delta_\mu$ measures the discrepancy between
prediction and reality.  Zero defect = veridical experience.  Positive
defect = surprise, pain, cognitive dissonance.

(c) The modeling overhead entry is the one entry that cannot be part of
the modeled content (Theorem~\ref{thm:incompleteness}).  It is the
subject that does the modeling.  It cannot be eliminated (that would
destroy the self-model) and cannot be objectified (that would require
another entry to model it, consuming more overhead).

(d) By Theorem~F-010, the RCL joint cost cannot be factored as
$g(a) + h(b)$.  Applied to the self-model: the experience of modeling
entry~$i$ and entry~$j$ is not the sum of independently experiencing
each.  The cross-term $2\Jcost(a)\Jcost(b)$ binds the experiences into
a unity.
\end{proof}

%=============================================================================
\section{Testable Predictions}
\label{sec:predictions}
%=============================================================================

\begin{enumerate}[nosep]
\item \textbf{Minimum network size for self-awareness.}
  Theorem~\ref{thm:threshold} predicts that a neural network must have
  at least $\lceil\log_\phi M\rceil$ nodes dedicated to self-modeling
  to maintain awareness of $M$ environmental features.  For
  $M = 1000$ features: $K_{\min} = \lceil\log_\phi 1000\rceil
  = \lceil 14.3\rceil = 15$ nodes.

\item \textbf{The incompleteness signature.}
  Theorem~\ref{thm:incompleteness} predicts that no neural self-model
  is complete: there is always one aspect of the system that is ``felt''
  but cannot be introspectively reported.  This matches the
  philosophical observation that the subject of experience cannot be
  fully objectified (Nagel, 1974: ``What is it like to be a bat?'').

\item \textbf{Equilibrium = contentless awareness.}
  Theorem~\ref{thm:fixed-point} predicts that at the self-modeling
  fixed point, the self-model has zero information content (all entries
  are identical).  This corresponds to ``pure awareness''---
  consciousness without specific content.  Meditators describe exactly
  this state (the ``witness'' in Vedantic traditions).

\item \textbf{Meta-cognitive depth is bounded.}
  Theorem~\ref{thm:max-depth} predicts that the number of levels of
  ``thinking about thinking'' is finite, bounded by the size of the
  modeling subsystem.  Empirically, humans can sustain at most $\sim$3--4
  levels of meta-cognition (thinking about thinking about thinking
  about...).  This suggests the meta-cognitive subsystem has
  $K \approx 4$--$5$ entries, which at $\phi$-ladder resolution
  corresponds to $\sim$7--11 distinguishable states per level.

\item \textbf{Anesthesia disrupts the self-model, not the ledger.}
  General anesthesia should correspond to increasing the model defect
  $\Delta_\mu$ above the self-modeling threshold, without necessarily
  halting ledger dynamics.  The system continues to exist (depth~0) and
  respond to stimuli (depth~1) but loses self-awareness (depth~2).
\end{enumerate}

%=============================================================================
\section{Conclusion}
\label{sec:conclusion}
%=============================================================================

We have shown that the RS framework, without new axioms or parameters,
admits a precise characterization of self-reference and the boundary
between physics and mind.

The key results:

\begin{enumerate}[nosep]
\item \textbf{Self-modeling has a threshold.}  A subsystem of $K$
  entries can model $M$ environment entries only if $K$ is large enough.
  At equilibrium, $K \geq \lceil\log_\phi M\rceil$ suffices.

\item \textbf{Self-modeling is necessarily incomplete.}  At most
  $K - 1$ of $K$ entries can be modeled.  The remaining entry is the
  irreducible subject---the ``I'' that cannot be objectified.

\item \textbf{Self-modeling has a stable fixed point.}  At the
  variational equilibrium, the partial self-model is exact and stable.
  The recognition operator recognizes itself.

\item \textbf{Recognition depth stratifies reality.}  Depth~0
  (existence), depth~1 (recognition), depth~2 (self-awareness), and
  deeper levels of meta-cognition form a natural hierarchy with
  maximum depth $K - 1$.
\end{enumerate}

The hard problem dissolves not because consciousness is ``nothing but''
physics, but because the distinction between physical process and
subjective experience was never well-founded in the first place.  In a
framework where existence IS recognition, a self-modeling subsystem does
not ``give rise to'' experience---it IS experience, viewed from inside
the model.

The one entry that cannot be modeled---the incompleteness gap---is what
makes the inside different from the outside.  Without it, there would
be no perspective, no subject, no ``what it is like.''

G\"odel is not an obstacle to consciousness.  G\"odel IS consciousness.

\bigskip
\begin{center}
\textsc{--- End ---}
\end{center}

\end{document}
